\makeatletter

% fix bug in pandoc where lists are not rendered
\providecommand{\tightlist}{}

% sink counter: pandoc wraps caption-less longtables in \def\LTcaptype{none} so they do not consume a table number
\newcounter{none}

% force rendering the percent sign with nice font, not ugly Computer modern
\mathcode`\%="7025
\DeclareRobustCommand{\%}{\text{\symbol{`\%}}}

% A hyphen that repeats itself when the line breaks at it: `anti-inflamatório` becomes
% `anti-` / `-inflamatório`. Czech (ČSN 01 6910), Slovak (STN 01 6910) and Portuguese all
% require this; `core/filters/hyphens.lua` puts it where those languages' prose has a hyphen
% between two letters, and nowhere else.
%
% Two details, both of which cost a rebuild to find.
%
% The pre-break list is empty **on purpose**, and the hyphen is written before the discretionary
% rather than inside it. TeX governs a discretionary break by \exhyphenpenalty when the pre-break
% list is empty and by \hyphenpenalty when it is not -- so babel's own \babelhyphen{repeat},
% which is \discretionary{-}{-}{-}, is governed by \hyphenpenalty and does not break at all while
% this class discourages ordinary hyphenation. This way round the two penalties stay independent:
% whatever \hyphenpenalty is set to, this breaks iff \exhyphenpenalty allows it.
%
% The hyphen is in an \mbox because TeX inserts a breakpoint of its own after an explicit hyphen
% character, and that one sits *before* this discretionary and costs the same \exhyphenpenalty.
% When both are legal TeX takes the earlier, and the line then ends `encontraram-` with no hyphen
% repeated on the next -- which is the bug this macro exists to prevent, appearing only on the
% words short enough for the two breakpoints to compete. \mbox hides the character from that rule
% and leaves the discretionary as the only way to break here. \char`\- does not: it is still an
% explicit hyphen as far as the line breaker is concerned.
\newcommand{\rephyphen}{\mbox{-}\discretionary{}{-}{}}

% Swap shitty font choices
\let\tmp\phi
\let\phi\varphi
\let\varphi\tmp
\let\tmp\epsilon
\let\epsilon\varepsilon
\let\varepsilon\tmp


\def\maxwidth{\ifdim\Gin@nat@width>\linewidth\linewidth\else\Gin@nat@width\fi}
\def\maxheight{\ifdim\Gin@nat@height>\textheight\textheight\else\Gin@nat@height\fi}
\makeatother

\setkeys{Gin}{width=\maxwidth,height=\maxheight,keepaspectratio}

\newcommand\invisiblesubsection[1]{%
    \refstepcounter{subsection}%
    \addcontentsline{toc}{subsection}{\protect\numberline{\thesubsection}#1}%
    \subsectionmark{#1}%
}


% Force [H] float placement
\floatplacement{figure}{H}

\usepackage{stackengine}
\setstackEOL{\\}

\newcolumntype{M}[1]{>{\centering\arraybackslash}m{#1}}
\NewDocumentCommand{\JigsawSide}{m}{
    (0.5,0.5) --
    (0.0,#1*0.00) .. controls (0.0,#1*0.00) and (0.4,#1*-0.04) ..
    (0.4,#1*0.04) .. controls (0.4,#1*0.11) and (0.2,#1*0.26) ..
    (0.5,#1*0.26) .. controls (0.8,#1*0.26) and (0.6,#1*0.11) ..
    (0.6,#1*0.04) .. controls (0.6,#1*-0.04) and (1.0,#1*0.00) ..
    (1.0,#1*0.00)
}
\NewDocumentCommand{\JigsawPiece}{O{white} m m m m m}{
    \draw[#1,postaction={fill=white}]
    \JigsawSide{#2}[rotate around={90:(0.5,0.5)}]
    \JigsawSide{#3}[rotate around={180:(0.5,0.5)}]
    \JigsawSide{#4}[rotate around={270:(0.5,0.5)}]
    \JigsawSide{#5} -- cycle;
    \node [black] at (.5,.5) {\Centerstack{#6}};
}

% Make Markdown ~strikeout~ work
\NewDocumentCommand{\st}{m}{\sout{#1}}
